% Instructions to modify this document:
% * Remember to ALWAYS execute "git pull" BEFORE any commit you make!
% * Use the \ToDo{...} command to remark tasks which still need to be done. Add your name in the comment.
% * Use the \input{file.tex} command to split the document into several parts
% * Do not change the current LaTeX coding style to yours. The style and format should be homogeneous along sections.

% To convert .dia diagrams into PDF:
% 1) Create the diagram with dia
% 2) Export it as .eps
% 3) use epstopdf to convert to PDF


\documentclass[a4paper,12pt]{article}

\usepackage[utf8]{inputenc}
\usepackage{amsmath,graphicx}
\usepackage{bm}
\usepackage{amssymb}
\usepackage{algorithm}
\usepackage{algpseudocode}
\usepackage{subfigure}
\usepackage{ifpdf}
\usepackage{url}
\usepackage{color}
\usepackage[hidelinks]{hyperref}
\usepackage{multirow}
\usepackage{datetime}
\usepackage{comment}
\usepackage{float} % To put figures in their exact place with \begin{figure}[H]
\usepackage{longtable}
\usepackage{tabularx}
\usepackage{listings}
\usepackage{xcolor}


% Definitions and commands
\def \np{\vskip 0.25 cm}
\def \ap{\vskip 0.15 cm}

\lstset{language=Bash, basicstyle=\color{gray}}

\newcommand{\ToDo}[1]{\textcolor{magenta}{\textbf{[ToDo]} \textbf{#1}}}
\newcommand{\miguel}[1]{\textcolor{magenta}{\textbf{[Miguel]} \textbf{#1}}}


\begin{document}

\begin{titlepage}

\begin{center}
\vspace*{-1in}

\vspace*{0.6in}
\begin{Large}
\textbf{The IPOL Demo System 2.0 \\Systeam Administration Documentation} \\
\end{Large}

\vspace*{0.6in}

\small{Compiled on \today\ at \currenttime}

\vspace*{0.6in}
\rule{80mm}{0.1mm}\\
\vspace*{0.1in}
\end{center}

\end{titlepage}

This document contains System Administration documentation for the IPOL Demo System 2.0.
\vspace*{0.6in}


%\maketitle
\newpage

\tableofcontents
\newpage
\listoffigures
\newpage



\section{System administration}

\subsection{Crontab scripts}
The system needs to periodically execute administration script, to create backups, clean up old data, or to perform continuous integration tests. This section describes these scripts.

To modify the crontab configuration, the command is {\tt crontab -e}. The format of the crontab lines is: minute (m), hour (h), day of month (dom), month (mon), and day of week (dow), or use '*' in these fields (for 'any').

\begin{itemize}
    \item \textbf{Backups}: they are done by this script: {\tt ipolDevel/sysadmin/backup.sh}
Since we do not want to the backup is not done just because some files can not be read because of their permission, we execute the backup script as the root user (use {\tt sudo crontab -e}).

We use the following line to run the backup script once per week at 3:00h:

    {\tt 0 3 * * 1 /home/ipol/ipolDevel/sysadmin/backup.sh}

    \item \textbf{Cleanup}: the script {\tt ipolDevel/sysadmin/cleanup.sh} does a cleanup of all temporary files older than one month.

    {\tt 0 5 * * 1 /home/ipol/ipolDevel/sysadmin/cleanup.sh}

    \item \textbf{PyLint report}: the script {\tt ipolDevel/ci\_tests/pylint.sh} uses PyLint to create code quality reports. It executes weekly in the integration server and sends an email to the addresses in the file {\tt send\_to.txt}

    {\tt 0 15 * * 2 /home/ipol/ipolDevel/ci\_tests/pylint.sh}

    \item \textbf{pdflatex report}: the script {\tt ipolDevel/ci\_tests/pdflatex.sh} checks if the Latex documentation of the project compiles and sends and email otherwise.

    {\tt 0 15 * * 2 sudo -u ipol /home/ipol/ipolDevel/ci\_tests/pdflatex.sh}

    \item \textbf{Reboot}: the script {\tt ipolDevel/sysadmin/on\_reboot.sh} starts all IPOL modules, nginx and the shared folder in the production machines, and sends a warning email. To execute it, the following line needs to be added with {\tt crontab -e} as root.

    {\tt @reboot ipolDevel/sysadmin/on\_reboot.sh}
\end{itemize}

\section{GitHub deploy key}
The IPOL GitHub repository might be private. Thus, is any server needs to obtain the complete source code of the demo system, a \emph{deploy key} is needed. The deploy key is added to the GitHub configuration and each server needs to get the sources using the deploy key.

For example:

\begin{verbatim}
ipol@smartalgo:~/ipolDevel$ cat ~/.ssh/config 
host github_deploy
      hostname github.com
      user git
      identityfile ~/.ssh/id_rsa_deploy
\end{verbatim}

\vspace{0.15cm}

\begin{verbatim}
ipol@smartalgo:~/ipolDevel$ cat .git/config 
[core]
	repositoryformatversion = 0
	filemode = true
	bare = false
	logallrefupdates = true
[remote "origin"]
	#url = https://github.com/mcolom/ipolDevel
	url = git@github_deploy:mcolom/ipolDevel.git
	fetch = +refs/heads/*:refs/remotes/origin/*
[branch "master"]
	remote = origin
	merge = refs/heads/master
\end{verbatim}

We use the host \emph{github\_deploy} in order to pick the deploy key in the SSH connection.

\section{Setup of a new machine}
When a new machine is added to the IPOL infrastructure, the following needs to be done:
- Create the {\tt ipol} user in a new machine: {\tt adduser --disabled-password ipol}

- Install bash-completion {\tt apt-get install bash-completion}

- Install cron-apt: {\tt apt-get install cron-apt}

And edit /etc/cron-apt/config:
\begin{verbatim}
APTCOMMAND=/usr/bin/apt-get
OPTIONS="-o quiet=1 -o Dir::Etc::SourceList=/etc/sources.list"
SYSLOGON="always"
\end{verbatim}

- Configuration of the email server to be able to send emails form any of the IPOL servers: configure with {\tt dpkg-reconfigure exim4-config} and be sure to choose the option to send email on the Internet, not local.

Install mail with {\tt apt-get install mailutils}

To send an email: {\tt cat email.txt | mail -s "This is the subject" john.doe@example.com,banania.guy@example.com}


\begin{comment}
- Setup priorities and memory limits to the processes of the users.

Added this to /etc/security/limits.conf:
\begin{verbatim}
*               soft     priority            17
ipol            soft     priority            0
\end{verbatim}

We need to set also limits for the memory usage per user. It'd be great if we could do it as a percentage of the free memory.

git@github_deploy:mcolom/ipolDevel.git
To limit the memory which a group of users can use, we should use {\tt cgroups} along with the {\tt Cgred} daemon.

More info:

\url{https://www.digitalocean.com/community/tutorials/how-to-limit-resources-using-cgroups-on-centos-6}

\url{https://access.redhat.com/documentation/en-US/Red_Hat_Enterprise_Linux/6/html/Resource_Management_Guide/sec-Moving_a_Process_to_a_Control_Group.html}
\end{comment}

To clone the ipolDevel repository in a new machine, the SSH access must be configured with the RSA deploy key so ssh can use it to clone the repository. This allows to clone the repository even if it eventually becomes private.
%
The deploy private key must be copied to {\tt .ssh/id\_rsa\_deploy} and the following lines need to be added to .ssh/config:

\begin{verbatim}
host github_deploy
      hostname github.com
      user git
      identityfile ~/.ssh/id_rsa_deploy
\end{verbatim}

Then, clone the repository with {\tt git clone git@github\_deploy:mcolom/ipolDevel.git}

\section{Nginx}

\subsection{Installation}

\begin{enumerate}
\item Debian

\begin{enumerate}
\item Uninstall previous nginx version
\begin{lstlisting}[language=Bash]
{prod} ~/ipolDevel$ apt-get remove nginx-common
\end{lstlisting}

\item Download the key from: \url{http://nginx.org/keys/nginx_signing.key} and execute the following command with it
\begin{lstlisting}[language=Bash]
{prod} ~/ipolDevel$ sudo apt-key add nginx_signing.key
\end{lstlisting}

\item Replace codename with Debian distribution codename, and append the following to the end of the /etc/apt/sources.list file:
\begin{enumerate}
\item deb http://nginx.org/packages/debian/ codename nginx
\item deb-src http://nginx.org/packages/debian/ codename nginx
\end{enumerate}

\item Update repositories
\begin{lstlisting}[language=Bash]
{prod} ~/ipolDevel$ apt-get update
\end{lstlisting}

\item Install nginx from the repository
\begin{lstlisting}[language=Bash]
{prod} ~/ipolDevel$ apt-get install nginx
\end{lstlisting}

\end{enumerate}
\item Ubuntu
\begin{enumerate}
\item Install nginx from the repository
\begin{lstlisting}[language=Bash]
{prod} ~/ipolDevel$ apt-get install nginx
\end{lstlisting}
\end{enumerate}
\end{enumerate}

\subsection{Configuration}

To configure the reverse proxy, we will use the sites-enabled, sites-available directories structure, which are located at:
\begin{lstlisting}[language=Bash]
/etc/nginx/sites-available
/etc/nginx/sites-enabled
\end{lstlisting}

\begin{enumerate}
\item Edit the config file /etc/nginx/nginx.conf, adding the following include at the end of the http subsection:
\begin{lstlisting}[language=Bash]
include /etc/nginx/sites-enabled/*;
\end{lstlisting}
\item Edit the site configuration file: sites-available/default, configuring the subsections:
\begin{enumerate}
\item Upstream
\begin{enumerate}
\item Upstream name. E.g.: ipol\_webapp\_server
\item Gunicorn socket file absolute location.
\end{enumerate}
\item Server
\begin{enumerate}
\item Listening port. The port to serve the site with static files. E.g.: 8000
\item Server name. E.g.: ipolcore.ipol.im
\item A resolver, necessary to set variables in the configuration for Debian distributions.
\begin{lstlisting}[language=Bash]
resolver 213.186.33.99;
\end{lstlisting}
\item Increase timeout for send and receive requests.
\begin{lstlisting}[language=Bash]
proxy_send_timeout	600;
proxy_read_timeout	600;
send_timeout		600;
\end{lstlisting}
\item Access log absolute location.
\item A location for each module.
\begin{lstlisting}[language=Bash]
location /api/<module>/ {
  rewrite ^/api/<module>/(.*) /$1 break;
  proxy_pass  http://$host:<port>;
}
\end{lstlisting}
\item Static files absolute location.
\begin{lstlisting}[language=Bash]
location /cp/static/ {
  alias  /home/<user>/IPOLWEBAPP_STATIC/;
}
\end{lstlisting}
\item Proxy pass, using upstream name configured previously.
E.g.:
\begin{lstlisting}[language=Bash]
proxy_pass http://ipol_webapp_server;
\end{lstlisting}

\end{enumerate}
\end{enumerate}
\item Enable the site by creating a symbolic link to default configuration:
\begin{lstlisting}[language=Bash]
ln -s /etc/nginx/sites-available/default /etc/nginx/sites-enabled/
\end{lstlisting}
\item Check Nginx configuration
\begin{lstlisting}[language=Bash]

\end{lstlisting}
\end{enumerate}

\subsection{Management}

\begin{itemize}
\item Start service:
\begin{lstlisting}[language=Bash]
{prod} ~/ipolDevel$ sudo nginx
\end{lstlisting}
\item Reload the configuration file:
\begin{lstlisting}[language=Bash]
{prod} ~/ipolDevel$ sudo nginx -s reload
\end{lstlisting}
\item Stop service:
\begin{lstlisting}[language=Bash]
{prod} ~/ipolDevel$ sudo nginx -s stop
\end{lstlisting}
\item Check configuration:
\begin{lstlisting}[language=Bash]
{prod} ~/ipolDevel$ sudo nginx -t
\end{lstlisting}
\end{itemize}


\section{Control Panel}
To operate in a production environment is necessary to install and configure some extra dependencies. In this case, Gunicorn and Nginx are 
required for this purpose. The first one will work as the application server, while the second is a reverse proxy. 

\begin{itemize}
\item \textbf{Reverse Proxy} is a type of proxy that retrieves resources on behalf of a client from one or more servers. 
These resources are then returned to the client as if they originated from the proxy server itself. The reverse proxy in this case is 
implemented with NGINX and is used to redirect requests from port 80 to the Control Panel listening port.
\item \textbf{WSGI} is a specification to connect web servers with web applications or frameworks for the Python programming language. 
In this case the Gunicorn is the responsible of this purpose.
\item \textbf{Web Application} is a is a client–server software application in which the client runs in a web browser. In the Control Panel this 
is implemented with Django.
\end{itemize}

\subsection{Communication between the Client and Control Panel}
The communication between the client and the Control Panel varies according to the environment in which it is performed.  
The differences between them are described below.

The communication between the client and the Control Panel in local is stablished by a reverse proxy. This reverse proxy is implemented with
NGINX and allows the client to do the requests in the 80 port. That request is intercepted by NGINX and send to the 8000
port where Django is running. The other purpose of NGINX is to serve the static content and remove that responsability from Django.
The diagram of this configuration is represented by Figure ~\ref{fig:local_architecture}.

\begin{figure}[!ht]
\centering
\includegraphics[width=0.7\columnwidth]{images/local}
\caption{The communication between the client and the Django web server with the reverse proxy in the local server.}
\label{fig:local_architecture}
\end{figure}

In the production and integration servers is where is located the most complete connection. As in the local server, a reverse proxy is needed to 
make the redirections to hide to the client in wich port is Django running and to provide all the static files. In addition, in the production 
and integration servers a WSGI is also needed to receive the http request and send them to Django. This WSGI is implemented with Gunicorn and 
communicates with the NGINX through a UNIX socket. The diagram of this configuration is represented by Figure ~\ref{fig:prod_int_architecture}.


\begin{figure}[!ht]
\centering
\includegraphics[width=1\columnwidth]{images/prod_int}
\caption{The communication between the client and the Django web server with the reverse proxy and the WSGI in the production and integration servers.}
\label{fig:prod_int_architecture}
\end{figure}


\subsection{Gunicorn: installation}

Install from console:
\begin{lstlisting}[language=Bash]
{prod} ~/ipolDevel$ pip install gunicorn
\end{lstlisting}

\subsection{Gunicorn: configuration}
In order to manage the gunicorn service we will use a script, which is in charge of:
\begin{itemize}
\item Open the virtual environment, if used. If not, comment the line with \#
\item Collect the static files.
\item Start the web server.
\end{itemize}

In this guide we will use the following path/name for the mentioned script:
\begin{lstlisting}[language=Bash]
ipolDevel/gunicorn_start.sh
\end{lstlisting}

Open the script and configure the parameters:
\begin{itemize}
\item NAME: Name of the application.
\item DJANGODIR: Django project directory path (absolute).
\item SOCKFILE: Gunicorn socket path (absolute).
\item VENV: Virtual environment name (if used).
\item USER, GROUP: User and group names, to run the server as. Both names must be valid at the current server.
\item NUM\_WORKERS: Gunicorn worker processes number.
\item DJANGO\_SETTINGS\_MODULE: Django settings name.
\item DJANGO\_WSGI\_MODULE: Django WSGI module name.
\item Gunicorn port number. In this guide we will use port 8001. Consider that this is NOT the port used to access the final site, because it will not serve the static files (DEBUG = False in settings.py). Another port will be configured in Nginx for that purpose.
\item Finally, add execution permissions to the script.
\end{itemize}

\subsection{Gunicorn: management}
\begin{itemize}
\item Start service by executing the script:
\begin{lstlisting}[language=Bash]
{prod} ~/ipolDevel$ ./gunicorn_start.sh
\end{lstlisting}

\item Stop service: Ctrl + C
\end{itemize}



\subsection{Accessing the Control Panel}
To access the Control Panel the user will need to access the specified URL, and login with the following credentials:

\begin{itemize}
\item URL: \url{http://ipolcore.ipol.im/cp/login/}

\item User: ipolcpadmin

\item Password: gy54g7x2 \miguel{Do we really want to write password in the public doc?}
\end{itemize}

\bibliographystyle{plain}
\bibliography{biblio}

\end{document}
% End of document
